\section*{Ringraziamenti} 
Eccoci qui. In un normale venerdì, dove tutto è finito. Nel momento che aspettavate tutti, quello che qualcuno definirebbe «instagrammabile», l'attimo in cui tutti tirano fuori il telefono per fare foto e video.
Oggi, proprio adesso, vorrei chiedervi di fermarvi un istante. Mettete via i telefoni: vorrei solo che mi ascoltaste. Perché in queste righe ho racchiuso, nel bene e nel male, tutto quello a cui ho pensato quando, arrivata alla fine della tesi, mi sono fermata e mi sono chiesta: «Ma in questi quattro anni, chi c'è stato davvero?»
\\
In primis ringrazio \textit{me stessa}. Ringrazio la me di diciotto anni che un giorno ha deciso: \textit{«Ora vado via, vado a Roma»}. Quella me coraggiosa e impaurita che ha chiuso una valigia, è salita su un aereo ed è atterrata a 900km di distanza. A quella me che ha iniziato a vivere per davvero, forse per la prima volta esattamente come desiderava, senza vincoli, restrizioni, e senza essere un bene da dividere.
Roma per me è stata libertà. Quando sono atterrata qui per la prima volta mi sono detta: \textit{«Gio, stai costruendo il tuo futuro. Qui non sei “la figlia di”, nessuno ti conosce, nessuno ti giudica».} Questa città per me ha il sapore di indipendenza, futuro, lavoro, sacrifici, obiettivi, lacrime, nuove amicizie ed esperienze. Una città che non si ferma mai, né di giorno né di notte, frenetica, dove tutti corrono per raggiungere qualcosa... anche se io non ho ancora capito bene cosa.
Poi però mi fermo, mi guardo indietro e vedo quella Giorgia di quattro anni fa: piccola, piena di sogni, per certi versi timida, partita da quel paesino che ha tanto amato e tanto odiato. Quel paese che profumava di salsedine, agrumi, della cucina di nonna, di calma, dialetto, famiglia e tradizioni. Ma anche di \textit{curtighiu}, di giudizi e sentenze.
Però devo dire che, nel bene e nel male, un po' mi manca, e mi è mancato.
Non sempre, ma in alcuni momenti più di tutti: quando durante la sessione estiva vedevo le amiche al mare, quando cucinavo da sola pensando a come a casa di nonna fossi una principessa, o quando aprivo il frigo deserto e toccava andare a fare il cosidetto spesone per la milionesima volta.
Perché è solo quando una cosa non ce l'hai più a portata di mano che ne capisci davvero il valore. 
Oggi guardo la me di adesso, mi giro e vedo la me di quattro anni fa, che da timida è diventata estroversa (che se poco poco lo ero già), là me che ha fatto amicizia, ha litigato, ha trovato il suo posticino in quel laboratorio, alla me che si è vissuta questa città sia di giorno che di notte, ma sopratutto un grazie va alla me che non ha mai mollato: ho sempre saputo di non essere stata mai un 10/10 su nulla, però ogni obiettivo che mi imponevo lo raggiungevo, con fatica, lacrime e lamenti perché diciamoci la verità: “mamma hai partorito un lamento costante, ma che non si arrende mai”.
\\\\
Ora passiamo a chi, in silenzio e borbottando, ha subito questa mia decisione. 
\\
\textit{Alla mia famiglia}, a mamma e papà, forse definire la nostra con il termine “famiglia” suona inusuale, ma per me lo siete stati e lo sarete per sempre. 
Se oggi qualcuno mi chiede “ma la tua famiglia?” io risponderò che è una delle cose parecchio complicate da poter descrivere, ma a modo suo, nel suo piccola è bella. 
Di certo, non è la famiglia del mulino bianco, gli spiegherò che è fatta da mamma da papà, da 2 case, da 2 natali…ma è bella pure così
Spero oggi, entrambi siate fieri di me, della donna che sono diventata, della figlia che sono stata in questi anni. Grazie ad entrambi, in modo diverso avete contribuito a formare la donna che sono diventata oggi.
\\
Grazie \textit{Mamma}, per avermi sostenuto sempre in tutto quello che mi passava per la mente, per ogni decisione che ho preso, tu eri li ad appoggiarmi sempre o la maggior parte delle volte. Grazie per non avermi mai fatto pensare che non avrei mai potuto fare quello che volevo o diventare ciò che volevo essere. Grazie per essere la mia più grande sostenitrice. Grazie, per aver capito che avevo bisogno di andare via, nonostante significava allontanarmi da te, per non avermi mai tarpato le ali.
Spero un giorno di essere la metà di quello che sei e di avere la metà della forza che hai avuto tu nella tua vita. Ricorda sempre, sii fiera di te stessa, non ti vergognare mai per come sei, perché con tutte le tue imperfezioni per me sei perfetta così. 
\\
Grazie \textit{Papà}, per avermi fatto capire chi voglio e chi non voglio essere. Tu instancabile lavoratore, mi hai insegnato la tenacia di non mollare quando gli altri non credono in te.
A volte mi fermo, penso che vorrei essere ancora quella bimba di 5 anni, che come diceva la canzone di Mia Martini, era innamorata del padre, quella bimba di cinque anni che vedeva in te il suo principe azzurro, quella che aspettava le domeniche alterne per fare una passeggiata a Marina e mangiare il gelato insieme. Crescendo, però, la favola del principe azzurro è svanita e oggi spesso ci scontriamo, perché guardiamo il mondo con lenti completamente diverse. A volte mi sento un po' la tua «figlia sgangherata», e ammetto che «ho lottato per cambiarti», anche se poi realizzo che «ci vorrebbe un'altra vita» proprio come dice la canzone, per metterci d'accordo su tutto e vedere il mondo in modo uguale.
Eppure, grazie per essere qui, in uno dei giorni più importanti della mia vita. Vederti lì per me significa tutto. Spero che tu oggi sia fiero della donna che sono diventata, perché dietro di lei c'è e ci sarà sempre quella bambina con cui andavi felice a prendere il gelato a Marina.
\\
\textit{Ai miei parenti}, nonni, zii, cugini. A chi c’è stato, a chi non c’è stato e a chi avrebbe voluto esserci ma non ha potuto. A chi, quando scendevo o tramite un messaggio, mi diceva: «Giò, come va a Roma? Come va l’università? Laureati così salgo su». A voi, alle nostre piccole famiglie che per me insieme formano una sola grande famiglia, voglio dire grazie. Grazie per l’infanzia vissuta insieme, se oggi sono la donna che vedete davanti a voi, lo devo interamente a come sono cresciuta lì con voi. Ai Natali, ai compleanni, alle feste, alle mangiate in campagna da nonno, alle giornate a Sampieri con le mie cugine e a tutti quei momenti che mi mancano incredibilmente. 
Sappiamo tutti che  ogni famiglia ha la propria matassa, più o meno ingarbugliata, e noi abbiamo la nostra, perdiamo anni a capire come possiamo scioglierle questa matassa, oggi qui a tutti vorrei chiedervi una cosa che a nessuno viene in mente, ma che è la cosa più semplice, ovvero tagliare quei fili, per ricominciare a viverci tutti insieme come un tempo.
\\\\
Grazie alle mie “persone medicina”, le mie 4 amichette bellissime. Grazie perché voi conoscete una Giorgia che gli altri non conoscono, c’è chi di voi mi conosce da 22 anni e chi mi conosce da 9 anni, ma voi a differenza di tutti conoscete la vera giorgia, quella timida quando non conosce, quella matta quando vuole fare “minchiate”, quella “curtighiara”, quella ansiosa, quella che ripeteva sempre “arà, mi siddia”, quella disordinata, quella che appena ha conosciuto un ragazzo a Roma ha chiamato subito voi e vi ha detto “ho conosciuto un biondino con gli occhi azzurri che faceva lo chef” e voi avete subito detto “gioo manda foto”.
\\
Alle sorelle che ho scelto, Bene e Sara.
A Bene, tutti hanno delle amiche di infanzia..ma tu per me sei stata come una sorella per 23 anni. È vero abbiamo litigato, fatto pace, ma tutto ciò che ti ho sempre detto, bello o brutto, era come se parlassi da sorella, non più grande non più piccola, ma alla pari. Non mi ricordo il nostro primo incontro, avevamo solo 10 e 11 mesi ed eravamo in quell’asilo nido di val Verde. Però ricordo tutto il resto, il periodo della scuola materna con la maestra Pippa e Tiziana, i cinque bellissimi anni di elementare con la maestra Fabia, e le medie seppur in classi separate la mattina, insieme i pomeriggi in cui facevamo strumento, ovviamente stesso strumento, quel flauto traverso che ho tanto odiato, le prime uscite i sabati sera, il mio primo shot bevuto a quella feste di fine anno al liceo.
Nonostante tutto, noi non ci siamo mai divise, sì abbiamo litigato, fatto pace, frequentato licei e università diverse, ma nel momento del bisogno c’eravamo l’una per l’altra, ad oggi spero che non ci “allontaneremo” più di così. Perché infondo noi possiamo pure litigare tutti i giorni, per opinioni e visioni diverse del mondo, ma spero che tu ti ricorda che a casa mia ci sarà sempre la porta aperta.
A Sara, la mia gemellina, tu sei arrivata dopo Bene e nonostante ciò per me sei sempre stata la mia sorellina più piccola. Per me sarai sempre la mia piccola Saretta, colei che ho visto nel pancione di Rosina, colei con cui cantavo Violetta e con cui mi alleavo per fare i dispetti a Bene, colei che ho visto crescere ed è diventata una bellissima donna, ma sopratutto fortissima. 
Grazie per esserci state, per voi ci sarò sempre come in questi anni ci siamo state a vicenda nei momenti difficili.
\\
Alle mie migliori amiche, Franca e Sonda,
Con voi ho passo i bellissimi 5 anni di liceo, gli anni dei primi amori, delle prime rivelazioni vere, lo studio, la disperazione per latino e filosofia. Grazieee per avermi sopportarto in questi 9 anni con tutte le mie lamentele, grazie per essermi state vicine e per avermi supportato quando vi ho detto “vado a studiare a Roma”, grazieee perché ogni volta che scendo sembra che il tempo non sia passato ma si sia solo fermato, graziee a voi per essere come sietee. Se qualcuno mi chiede di voi? Io risposndo semplicemene “sono Franca e Sonda, il mio 10 su 10 su tutto”, perché diciamoci la verità io ho sempre detto e continuo a ripetere una cosa “mi sono scelta e circondata di persone migliori di me” voi non eravate solo le migliori della classe, non eravate solo migliori nello studio, eravate le migliori nella vita…e questo si è rivelato con il tempo. 
\\
A Sonda, conosciuta da tutti come Alessia.
Troppo complicato signorina Marinero spiegare a tutti perché la chiamo così, quindi mi concentrerò su altro. A te, che mi hai sempre detto: «Gio, tu che riesci a far predominare l’apollineo sul dionisiaco».
Grazie per avermi vista come neanche io riesco a vedermi. Non so davvero se io sia l’apollineo: la calma, la misura, l’ordine; o il dionisiaco: l'istinto puro, l'emozione e il caos. Mi sento dire da te che trasmetto tranquillità, eppure spesso devo ricordarmi anche io di respirare e imparare a gestire l’ansia.
Credo che tu mi abbia sempre descritta migliore di quanto io credessi di poter essere. Grazie perché, anche se non so se Donzy sia fiero di me, so per certo che lo sei tu. E la verità è che ci incastriamo perfettamente: se io sono al 60% apollinea, tu sei quel 40% dionisiaco che mi accende; e quando il mio caos dionisiaco prende il sopravvento in me, tu diventi esattamente la parte apollinea di cui ho bisogno per ritrovare la calma.
\\
A Fra, 

\\\\
Ed ora, a coloro che ho sconoscito in questo percorso.
\\
A Mario,
\\
Un grazie speciale va a Rosaria e Rodolfo, che mi hanno accolta come una figlia, a Claudio e Fede che mi hanno accolta come una sorella. A Ines, Franco, Alda, Giovi, Barbi, Paola e tutti gli altri che mi hanno fatto sentire parte di una famiglia.
\\\\
Ai miei migliori amici, 
A Chri, le parole su di te non bastano. Penso davvero che tu sia l’amico che ogni ragazza vorrebbe avere al suo fianco. Mi hai sempre ascoltata senza mai giudicarmi; al massimo, davanti ai miei racconti su delle persone, la tua risposta era: "pivello, ah gio io lo so fa meglio". Sei stato il mio confidente per le serate, le persone, le materie dell'università e i momenti difficili a casa. Insomma, tu di me sai tutto, e sarà sempre così.
A Pietro, grazie a te ho imparato il valore della tranquillità, del vero chill. Davanti a ogni mio sfogo la tua risposta era sempre "tranquilla gio, stai chill, si risolve", e alla fine avevi sempre ragione. Grazie per esserci stato, anche in quei momenti che per te, forse, tanto chill non erano. Grazie per quei caffè ai tavolini fuori del bar, dove tu puntualmente fumavi un numero k di sigarette, con k >= 1 (e k palesemente tendente a infinito). Con te ho imparato a bere bene, e che il mio corpo non regge tantissimo ma abbastanza. Grazie per avermi fatto vivere la vera Roma notturna: le serate al Largo Venue, il nostro amico bangla al Pigneto, il Nice e il Forte Prenestino. Quando avrai un momento di sconforto, ricordati di sorridere pensando a: "viva viva punk a bestia".
\\\\
Alle mie Coinqui,
A Fra e Anto o Anto e Fra, tanto cambiando l'ordine degli addendi il risultato non cambia. Per me siete state molto più di due semplici coinquiline con cui condividere una casa: siete state una boccata d'aria fresca. Siete quel genere di persone che, quando ti entrano nella vita, ti fanno dire: "Ma perché non le ho incontrate prima?". E io, in questi due anni, me lo sono chiesta spesso. Siete state la mia salvezza nei giorni "no"…quando in quei giorni, mi prendevate, puntualmente andavamo a finire nell’armadio di Fra che mi prestava i vestiti, Anto che mi truccava, e poi via... un pub, del buon alcol, un po' di musica e all'improvviso tutto passava. Siete state la leggerezza perfetta per i miei momenti più pesanti.
\\E poi, dopo un po', nella nostra folle quotidianità è arrivata Rosa. Mi piace pensare che non sia stato un caso, ma puro destino. Arrivata timida e introversa...e che ora, piano piano, stiamo irrimediabilmente rovinando! Sei diventata la mia traduttrice ufficiale quando le persone non mi capiscono, e l'unica che riesce a seguirmi quando mi lamento di qualche matematico o informatico che scrive teoremi indecifrabili, o quando me ne esco dicendo che so calcolare i percorsi meglio di Dijkstra (e altre cavolate simili!)
Non so quando lasceremo quella casa e non so come sarà vivere in una casa senza di voi. Ma anche se la vita adesso dovesse portarci altrove, voglio ricordarci così: tutte insieme davanti a un piatto cucinato da me, accompagnate da una canzone cantata a squarciagola da Anto, mentre Fra e Rosa (per fortuna) puliscono il campo di battaglia che ho lasciato ai fornelli.
\\\\
Alla Franci, all’amica che ho incontrato da adulta, a te che mi sei sempre stata vicina in questi tre anni, senza invidia e senza giudicarmi…sei un anima buona. Spero che fuori da qui, possa il nostro rapporto continuare, come quei weekend fuori da sole all’ikea o in un centro commerciale a parlare della qualsiasi cosa.
\\\\
A Compss, Francesco, tu sei stato la prima persona con cui ho scambiato mezza parola tra queste mura. Tutto è iniziato con un tuo messaggio indimenticabile: «Gio scusami, io sto male, però ti lascio il numero di uno dei miei amici, vi raggiungo dopo?». Chi l'avrebbe mai detto che quel mal di pancia avrebbe fatto nascere il nostro primo gruppo e mi avrebbe dimostrato che timida, in realtà, non lo sono mai stata.
Grazie per i pranzi a casa mia, per il gossip terapeutico, per il supporto coi regali di Natale, per ogni uscita e per la passione che ci accomuna per Scuderia. Grazie per ciò che è stato e per tutto quello che deve ancora arrivare. Ma non dimenticare mai: “Charles batte Hamilton su tutti i fronti”.
\\\\
E infine, la vera ciliegina sulla torta. L’ultimo ringraziamento, che non è ultimo per importanza, anzi tutto il contrario, va 
Al LAB25A, quel posto da cui qualunque persona sana di mente fuggirebbe a gambe levate dopo aver chiamato le forze dell'ordine e un sacerdote. Un giorno qualcuno di saggio mi disse: «Questo è il regno delle scimmie». La verità è che, una volta conosciute quelle scimmie, prima le ho criticate e poi, piano piano, ho imparato a volergli bene.
Grazie perché, tra alti e bassi, mi avete aiutata a capire chi sono e, soprattutto, chi non voglio essere. Grazie perché in questi quattro anni quel laboratorio è stato casa mia più della mia vera casa: ci ho fatto colazione, pranzato, cenato e, a volte, passato tutta la notte (diciamo non proprio dormendo). A quel posto che mi ha vista ubriaca durante gli hackathon, e versione mamma quando cercavo di rincorrere i piccini che giocavano a pallone facendosi richiamare dalla vigilanza, aiutando chi smarriva le chiavi e a tanto altro.
E ora? Non resta che ringraziare di cuore chi ha reso tutto questo indimenticabile: Franco, il Dottor Maestro, Pietro, Christian, Fra, Ivan, Mirco, Alex, Flavia, Davide, Dimu, Comps, Lorenzino, Vitto, Matteo, Luca, Masso, Alessia, Claudia, Tony $&$ i minions, Ilary e Francesca.